\begin{textsample}{POS Dim 3 – am – Score 24.00 – am/am\_inf\_27.txt}  \label{ex:f3_pos_040}
13.
As transições dos \textbf{bebês} e \textbf{criança} em seu processo educativo Pensar no fazer dessa etapa tão recente na história da educação implica pensar nos períodos de transição do bebê/criança em seu processo educativo, a fim de organizarmos o tempo, espaço e materiais de forma a integrar as etapas vividas, que passam por passagens e rupturas em seu processo de adaptação e desenvolvimento.
O primeiro deles já ocorre no \textbf{nascimento}, em que, pela primeira vez, o \textbf{bebê} precisar enfrentar o mundo externo.
É no contato com a família que o \textbf{bebê} vai se relacionando com as vivências de seu grupo.
A espécie humana é aquela que mais leva tempo para se tornar independente, mas o que poderia ser uma desvantagem na natureza constitui o meio pelo qual apreendemos a cultura e nos humanizamos.
Superada a ideia assistencial de \textbf{cuidado} de \textbf{bebês} e \textbf{crianças} de mães trabalhadoras, temos hoje a garantia do direito da \textbf{criança} à educação e o desafio da promoção da integração de \textbf{cuidar} e educar em espaços e \textbf{instituições} não \textbf{domésticas} públicas ou privadas, em períodos parciais ( mínimo de 4h ) e integrais ( mínimo de 7h ), no período diurno, com a garantia de profissional com formação adequada para esse atendimento.
Esse \textbf{bebê} que sai de casa para a \textbf{creche} vivencia uma importante mudança na sua \textbf{rotina} e no seu desenvolvimento.
Pessoas e ambientes diferentes lhes são apresentados.
Aqui, nossa preocupação deve estar na garantia do bem-estar emocional, planejando práticas que facilitem a transição e a adaptação ao novo contexto social.
É um processo novo e desafiador para os \textbf{bebês}, suas as famílias e também para as \textbf{instituições}, que devem promover um ambiente seguro e desenvolvente para esse \textbf{bebê} que está vivenciando a cada \textbf{dia}, novas experiências. \textbf{Cuidado} com o sono e com os objetos de transição podem tornar a experiência do \textbf{bebê} mais \textbf{agradável}.
Devemos permitir que a família traga para a \textbf{creche} objetos de acalento, que transmitem a sensação de \textbf{conforto} e segurança, como chupeta, fralda, bichinho de pelúcia etc., pois estes ajudam o \textbf{bebê} a ter um sono mais tranquilo e reduz o estresse da separação de casa.
A construção de uma \textbf{rotina} pensada a partir das interações e \textbf{brincadeiras} e com o \textbf{cuidar} e o educar pensado de forma indissociável é a base de uma proposta \textbf{diária} que cria \textbf{laços} de confiança dos \textbf{bebês}, \textbf{crianças} bem \textbf{pequenas}, \textbf{crianças} \textbf{pequenas} e suas famílias com a proposta de trabalho da \textbf{instituição}.
A \textbf{rotina} vai dizendo a sequência das coisas que estão por vir até chegar o momento de voltar para casa.
Conforme os \textbf{bebês} e as \textbf{crianças} vão percebendo esse estabelecimento de ações que se desencadeiam ao longo dos \textbf{dias} e semanas, a relação deles com o ambiente fica mais tranquila, à medida que a sensação de segurança e confiança são criadas com as pessoas e lugares novos.
A transição da \textbf{creche} para a \textbf{pré-escola} geralmente é marcada pela ideia de que haverá mais “ tarefas ” e que educar será mais preponderante que o \textbf{cuidar}.
Novamente, a indissociação destes princípios que regem o fazer na Educação \textbf{Infantil} deve prevalecer.
Em toda a etapa, estas ações estão vinculadas ao processo educativo da primeira \textbf{infância} e por isso, não deve haver professores e cuidadores com tarefas separadas, onde um ensina e, o outro \textbf{cuida} do corpo.
A \textbf{criança} que frequenta a \textbf{creche} chega à \textbf{pré-escola} com vivências de grupo e familiarizada com a \textbf{rotina} dos espaços educativos, o que geralmente contribui para uma adaptação mais tranquila.
Sabemos que isto não é regra e o \textbf{cuidado} com a ruptura da \textbf{instituição} \textbf{creche} ou casa para a entrada na \textbf{pré-escola} deve ser considerado.
A \textbf{criança} continua na mesma etapa e os eixos do currículo são os mesmos : interações e \textbf{brincadeiras}.
Superada a ideia de fase preparatória para o Ensino Fundamental, a \textbf{pré-escola} não é o espaço de aulas pré-alfabetizatórias que em nada contribuem para o processo de construção da compreensão dos usos sociais da escrita e da leitura.
A \textbf{brincadeira} de papéis e as atividades \textbf{plásticas} constituem a via pela qual as \textbf{crianças} mais se desenvolvem psiquicamente, e estas atividades devem ser privilegiadas nas turmas de \textbf{crianças} de 4 e 5 anos.
As diferentes linguagens devem ser diariamente exploradas na ampliação dos repertórios e na criação de atividades cada vez mais enriquecedoras de experiências.
Os direitos de aprendizagem devem ser garantidos em todas as etapas da Educação \textbf{Infantil}.
Segundo a BNCC, na transição da Educação \textbf{Infantil} para o Ensino Fundamental, precisamos garantir a “ integração e continuidade dos processos de aprendizagens das \textbf{crianças}, respeitando suas singularidades e as diferentes relações que elas estabelecem com os conhecimentos, assim como a natureza das mediações de cada etapa ” ( BRASIL, 2017, p. 51 ).
É urgente e necessário um diálogo entre as \textbf{instituições} de Pré-Escola e dos Anos Iniciais do Ensino Fundamental.
A ruptura vivida pela \textbf{criança} nessa transição é ressaltada pela organização dos espaços, \textbf{mobiliários} e não raras vezes, na postura dos profissionais que querem eliminar características da \textbf{infância}, como a \textbf{brincadeira}.
Lembremos que a \textbf{criança} de 6 anos ainda tem como atividade principal de aprendizagem e desenvolvimento a \textbf{brincadeira} de papéis e as atividades \textbf{plásticas}.
A escrita e leitura precisam continuar privilegiando os variados gêneros textuais, sempre trabalhados em contextos reais que criarão necessidade na \textbf{criança} pelo ato de ler e escrever.
As Diretrizes Curriculares Nacionais para a Educação \textbf{Infantil} ( BRASIL, 2009b ), asseguram que : Art. 10.
As \textbf{instituições} de Educação \textbf{Infantil} devem criar procedimentos para \textbf{acompanhamento} do trabalho pedagógico e para avaliação do desenvolvimento das \textbf{crianças}, sem objetivo de seleção, promoção ou classificação, garantindo : III – a continuidade dos processos de aprendizagens por meio da criação de estratégias adequadas aos diferentes momentos de transição vividos pela \textbf{criança} ( transição casa/instituição de Educação \textbf{Infantil}, transições no interior da \textbf{instituição}, transição creche/pré-escola e transição pré-escola/Ensino Fundamental ) ; Art. 11.
Na transição para o Ensino Fundamental a proposta pedagógica deve prever formas para garantir a continuidade no processo de aprendizagem e desenvolvimento das \textbf{crianças}, respeitando as especificidades \textbf{etárias}, sem antecipação de conteúdos que serão trabalhados no Ensino Fundamental.
Para realizar essa transição de forma sequencial e articulada, \textbf{instituições} de Educação \textbf{Infantil} e Ensino Fundamental devem ter uma organização, via secretaria de educação, que garanta continuidade nas experiências vividas pela \textbf{criança} nos primeiros anos escolares, que garanta, por meio da formação continuada, uma discussão que esclareça aos professores os anseios e necessidades de cada etapa e a forma pela qual a \textbf{criança} mais aprende e se desenvolve, sem promover antecipação de conteúdos que objetivam preparar para etapas posteriores, considerando os interesses e as necessidades das \textbf{crianças} no presente ; promover encontros entre os professores de \textbf{crianças} de 5 e 6 anos ; promover práticas que considerem o desenvolvimento integral dessas \textbf{crianças} e estreitar o diálogo entre a \textbf{instituição} de Pré-Escolar e as escolas de Ensino Fundamental próximas, promovendo visitas e encontros entre professores e \textbf{crianças}.
Essas são algumas formas possíveis de organizar uma passagem sem tensões e rupturas.
Como professores dos primeiros anos da \textbf{criança} no espaço formal de educação, precisamos defender arduamente o não encurtamento da \textbf{infância}.
Antecipar práticas do Ensino Fundamental e substituir o tempo e espaço das interações e \textbf{brincadeiras} negam os direitos das \textbf{crianças} de viver sua \textbf{infância} e ter a \textbf{brincadeira} como um elemento estruturante da sua cultura.
Estas práticas também enfadam e aborrecem a \textbf{criança} com imposições que ela ainda não vê sentido.
Quando elas sentirem a necessidade pela leitura e escrita e tiverem vivido as experiências que sustentam o aprendizado da escrita e da leitura, tendo sua função simbólica desenvolvida a cada \textbf{dia} por meio das ações \textbf{lúdicas} promovidas na Educação \textbf{Infantil}, elas escreverão e lerão com sentido e conscientes da função dessas linguagens.

% matched lemmas: acompanhamento, agradável, bebê, brincadeira, conforto, creche, criança, cuidado, cuidar, dia, diário, doméstico, etário, infantil, infância, instituição, laço, lúdico, mobiliário, nascimento, pequeno, plástico, pré-escola, rotina
\end{textsample}
